\begin{definition-proposition}
	\index{fenliecha@分裂叉 (split fork)}
	考虑任意范畴 $\mathcal{C}$ 中的图表
	\begin{equation*}\label{eqn:split-coeq}\begin{tikzcd}[column sep=large]
		A \arrow[shift left, r, "u"] \arrow[shift right, r, "v" description] & B \arrow[r, "h"] \arrow[dashed, l, bend left, "t"] & Z \arrow[dashed, l, bend left, "s"]
	\end{tikzcd}\end{equation*}
	使得实线部分交换, 亦即 $hu = hv$, 而且 $h$ 和 $v$ 各自有虚线所示的截面 (亦即右逆) $s$ 和 $t$, 满足 $sh = ut \in \End(B)$. 具此性质的实线部分图表称为\emph{分裂叉}. 它自动给出 $u$ 和 $v$ 在 $\mathcal{C}$ 中的余等化子.
\end{definition-proposition}
\begin{proof}
	考虑以下场景
	\[\begin{tikzcd}[row sep=small]
		& & W \\
		A \arrow[shift left, r, "u"] \arrow[shift right, r, "v"'] & B \arrow[r, "h"'] \arrow[ru, "k"] & Z
	\end{tikzcd} \qquad \begin{array}{ll}
		W \in \Obj(\mathcal{C}) \\
		ku = kv.
	\end{array} \]
	若存在 $\varphi: Z \to W$ 使得 $\varphi h = k$, 则 $\varphi = \varphi hs = ks$; 反之, $\varphi := ks$ 确实满足 $\varphi h = ksh = kut = kvt = k$. 这就验证了泛性质.
\end{proof}

不同于一般的余等化子, 分裂叉是``绝对''的: 它们对任意函子的像仍是分裂叉.

\begin{example}\label{eg:T-split-fork}
	设 $(T, \mu, \eta)$ 是 $\mathcal{C}$ 上的单子, 则任何 $(M, a) \in \Obj(\mathcal{C}^T)$ 都给出 $\mathcal{C}$ 中的分裂叉
	\begin{equation*}\begin{tikzcd}[column sep=large]
			T^2(M) \arrow[shift left, r, "Ta"] \arrow[shift right, r, "\mu_M" description] & T(M) \arrow[r, "a"] \arrow[dashed, bend left, l, "\eta_{TM}"] & M \arrow[dashed, bend left, l, "\eta_M"]
		\end{tikzcd} \quad \begin{array}{ll}
			a (Ta) = a \mu_M , & \eta_M a = (Ta) \eta_{TM}, \\
			a \eta_M = \identity_M , & \mu_M \eta_{TM} = \identity_{TM}.
	\end{array}\end{equation*}
	右侧列出的等式是定义 \ref{def:monad} 和 \ref{def:T-Alg} 的内容; 例如 $\eta_M a = (Ta) \eta_{TM}$ 缘于 $\eta: \identity_{\mathcal{C}} \to T$ 的自然性. 特别地, 上图将 $M$ 实现为 $Ta$ 和 $\mu_M$ 的余等化子.
\end{example}

\begin{definition}\label{def:Beck-split-pair}
	考虑函子 $G: \mathcal{D}\to \mathcal{C}$ 和 $\mathcal{D}$ 的一对态射 $f, g: X \rightrightarrows Y$. 如果存在态射 $h: G(Y) \to Z$ 使得图表
	\begin{equation*}\begin{tikzcd}
			G(X) \arrow[shift left, r, "Gf"] \arrow[shift right, r, "Gg"'] & G(Y) \arrow[r, "h"] & Z
	\end{tikzcd}\end{equation*}
	是 $\mathcal{C}$ 中的分裂叉, 则称 $(f, g)$ 是 $\mathcal{D}$ 中的 $G$-分裂对.
\end{definition}

设 $(f, g)$ 是 $G$-分裂对. 既然 $(Gf, Gg)$ 有余等化子, 一个自然的问题是能否将之提升为 $(f, g)$ 在 $\mathcal{D}$ 中的等化子, 以及此提升是否唯一. 余等化子是 $\varinjlim$ 的特例, 定义 \ref{def:create-limit} 关于函子生 $\varinjlim$ 或保 $\varinjlim$ 的概念为此提供了一套方便的术语.

留意到忘却函子 $U^T: \mathcal{C}^T \to \mathcal{C}$ 是定义 \ref{def:conservative} 所谓的保守函子.

\begin{lemma}\label{prop:UT-split-fork}
	设 $(T, \mu, \eta)$ 是 $\mathcal{C}$ 上的单子, 则函子 $U^T$ 生 $U^T$-分裂对的余等化子.
\end{lemma}
\begin{proof}
	设有 $\mathcal{C}^T$ 中的一对态射 $u, v: (M, a) \rightrightarrows (M', a')$ 和 $\mathcal{C}$ 中的分裂叉
	\begin{equation*}\begin{tikzcd}[column sep=large]
		M \arrow[shift left, r, "u"] \arrow[shift right, r, "v" description] & M' \arrow[r, "h"] \arrow[dashed, l, bend left, "t"] & M'' \arrow[dashed, l, bend left, "s"].
	\end{tikzcd}\end{equation*}
	由于上图给出 $\mathcal{C}$ 中的等化子, 问题归结为将 $M''$ 唯一地扩充为 $(M'', a'') \in \Obj(\mathcal{C}^T)$, 使得 $h$ 是 $\mathcal{C}^T$ 中的态射, 然后说明这给出 $u$ 和 $v$ 在 $\mathcal{C}^T$ 中的余等化子.
	
	分裂叉对任意函子的像仍是分裂叉, 故在实线部分的交换图表
	\[\begin{tikzcd}
		T(M) \arrow[shift left, r, "Tu"] \arrow[shift right, r, "Tv"'] \arrow[d, "a"'] & T(M') \arrow[r, "Th"] \arrow[d, "{a'}"] & T(M'') \arrow[dashed, d, "{a''}"] \\
		M \arrow[shift left, r, "u"] \arrow[shift right, r, "v"'] & M' \arrow[r, "h"'] & M''
	\end{tikzcd}\]
	中, 两行都是 $\mathcal{C}$ 中的余等化子, 从而泛性质给出唯一的 $a''$ 使得全图交换. 如能说明 $(M'', a'')$ 是 $T$-代数, 则上图右部交换相当于说 $h$ 是 $\mathcal{C}^T$ 中的态射.
	
	为此, 问题在于验证以下交换图表
	\[\begin{tikzcd}
		M'' \arrow[r, "{\eta_{M''}}"] \arrow[equal, rd] & T(M'') \arrow[d, "{a''}"] \\
		& T(M'')
	\end{tikzcd} \quad
	\begin{tikzcd}
		T^2(M'') \arrow[r, "{\mu_{M''}}"] \arrow[d, "{Ta''}"'] & T(M'') \arrow[d, "{a''}"] \\
		T(M'') \arrow[r, "{a''}"'] & M'' .
	\end{tikzcd}\]
	余等化子图表说明 $h$ 和 $Th$ 满 (同理, $T^2 h$ 满), 故易从 $M$ 和 $M'$ 的相应交换图表说明上图确实交换.
	
	最后来说明 $h$ 给出 $u$ 和 $v$ 在 $\mathcal{C}^T$ 中的余等化子. 设有 $\mathcal{C}^T$ 的态射 $k: (M'', a'') \to (W, b)$ 满足 $ku = kv$, 则在 $\mathcal{C}$ 中存在唯一态射 $\varphi: M'' \to W$ 使得 $\varphi h = k$. 问题归结为证明 $\varphi$ 实则是 $\mathcal{C}^T$ 的态射. 考虑图表:
	\[\begin{tikzcd}
		T(M') \arrow[r, "Th"] \arrow[d, "{a'}"'] & T(M'') \arrow[r, "T\varphi"] \arrow[d, "{a''}"] & T(W) \arrow[d, "b"] \\
		M' \arrow[r, "h"'] & M'' \arrow[r, "\varphi"'] & W
	\end{tikzcd}\]
	上下两行分别合成为 $Tk$ 和 $k$, 因为 $h$ 和 $k$ 是 $\mathcal{C}^T$ 中的态射, 图表的整个外框和左方块皆交换, 故右侧方块和 $Th$ 合成以后交换. 已知 $Th$ 满, 故右侧方块本身交换, 这就说明了 $\varphi$ 是 $\mathcal{C}^T$ 中的态射.
\end{proof}

\begin{lemma}\label{prop:Beck-FG-coeq}
	设函子 $G: \mathcal{D} \to \mathcal{C}$ 有左伴随 $F$, 而且 $G$ 对所有 $G$-分裂对 $f, g: X \rightrightarrows Y$ (定义 \ref{def:Beck-split-pair}) 都生相应的余等化子, 则下图给出余等化子:
	\begin{equation*}\begin{tikzcd}[column sep=large]
		FGFG(N) \arrow[shift left, r, "FG\varepsilon_N"] \arrow[shift right, r, "\varepsilon_{FG(N)}"'] & FG(N) \arrow[r, "\varepsilon_N"] & N,
	\end{tikzcd}\end{equation*}
	其中 $N \in \Obj(\mathcal{D})$.
\end{lemma}
\begin{proof}
	对图表取 $G$ 给出例 \ref{eg:T-split-fork} 在 $(M, a) := (G(N), G\varepsilon_N) = \mathbb{K}(N)$ 情形的分裂叉, 而 $G$ 生相应的余等化子. 故原图是余等化子.
\end{proof}
